\chapter{State of the Art}
\label{ch:sota}

Having outlined in the introductory chapter the application context of cooperative driving and platooning, this chapter explores the state of the art on the aspects that are directly relevant to the experimental work: platoon organization, longitudinal controllers, V2X technologies for intra-platoon communications, and multi-channel architectures for degradation management.

A cross-cutting theme in all the literature considered is the strong interdependence between networking and control. Communication quality cannot be evaluated as an isolated parameter, because it directly affects safety gaps and the platoon's ability to absorb disturbances without amplifying them \cite{dressler2019cooperative,giordano2019joint,segata2016safe}. For this reason, more recent approaches tend to favor multi-technology solutions and controlled degradation mechanisms rather than relying on a single interface or on overly abrupt fallbacks \cite{ishihara2015improving,segata2016platooning,segata2023multi-technology}.

The analysis concludes with a discussion of the \emph{SafeSwitch} mechanism proposed by Segata \emph{et al.} \cite{segata2023multi-technology}, which serves as the direct reference for the experimental part of the thesis. In this way, the chapter builds the bridge between the general framing already introduced in the introduction and the implementation described later on.
